% Part: first-order-logic
% Chapter: proof-systems
% Section: axiomatic-deduction

\documentclass[../../../include/open-logic-section]{subfiles}

\begin{document}

\iftag{FOL}
      {\olfileid{fol}{prf}{axd}}
      {\olfileid{pl}{prf}{axd}}

\olsection{بديهي \usetoken{P}{derivation}}

بديهي !!{derivation}s د !!{derivation} تر ټولو زاړۀ او ساده منطقي
نظامونه دي۔ د هغو !!{derivation}s يوازې د !!{sentence}s لړۍ دي۔ د
!!{sentence}s يوه لړۍ هله سم !!{derivation} ګڼل کېږي چې په هغې کښې
هره !!{sentence}~$!A$ له لاندې شرطونو څخه يو پوره کړي:
\begin{enumerate}
\item $!A$ بديهي اصل وي؛ يا
\item $!A$ د !!{sentence}s د ورکړل شوي سټ~$\Gamma$ يو !!a{element}
  وي؛ يا
\item $!A$ د استنتاج په کومه قاعده توجيه شوے وي۔
\end{enumerate}
د بديهي اصل کېدو دپاره $!A$ بايد د يو شمېر ټاکلو !!{sentence}
شېماوو له يوې سره بڼه ولري۔ د بديهي اصلونو د شېماوو ډېر داسې سټونه
شته چې د لومړۍ درجې منطق دپاره د !!{derivation} د قناعت وړ، سم او
بشپړ نظام ورکوي۔ ځينې د هغو نښلوونکو له مخې ترتيب شوي وي چې واک پرې
چلوي؛ مثلاً دا شېماوې
\[
!A \lif (!B \lif !A) \qquad !B \lif (!B \lor !C) \qquad (!B \land !C) \lif !B
\]
هغه عام بديهي اصول دي چې پر $\lif$، $\lor$ او~$\land$ واک چلوي۔ د
بديهي اصولو ځينې نظامونه غواړي د اصولو شمېر تر ټولو لږ وي۔ دا چې
کوم نښلوونکي بنسټيز وټاکل شي، ان داسې نظامونه هم موندل کېدے شي چې
يوازې له يوه بديهي اصل څخه جوړ وي۔

د استنتاج قاعده يوه شرطي وينا ده چې په !!a{derivation} کښې د
!!a{sentence} د توجيه کېدو دپاره بس شرط ورکوي۔ مودوس پوننس يوه ډېره
عامه داسې قاعده ده: وايي که $!A$ او $!A \lif !B$ له مخکښې توجيه
شوي وي، نو $!B$ توجيه شوے دے۔ دا معنٰی لري چې په !!a{derivation}
کښې هغه کرښه چې !!{sentence}~$!B$ لري، په دې شرط توجيه شوې ده چې
$!A$ او $!A \lif !B$ دواړه، د کوم !!{sentence}~$!A$ دپاره، په
!!{derivation} کښې له~$!B$ څخه مخکښې راغلي وي۔

پر بديهي !!{derivation}s ولاړه $\Proves$ اړيکه داسې تعريفېږي:
$\Gamma \Proves !A$ هله او يوازې هله چې داسې !!a{derivation} موجود
وي چې !!{sentence}~$!A$ ئې وروستے فارمول وي، او $\Gamma$ په هغه
!!{derivation} کښې د هغو !!{sentence}s سټ ونيول شي چې د پورته (۲) له
مخې توجيه شوي وي۔ $!A$ هله قضيه ده چې $!A$ داسې !!a{derivation}
ولري چې~$\Gamma$ پکې تش وي؛ يعنې په !!{derivation} کښې هره
!!{sentence} يا د (۱) له مخې توجيه شوې وي يا د~(۳) له مخې۔ مثلاً دا
!!a{derivation} ښيي چې $\Proves !A  \lif (!B \lif (!B \lor !A))$:
\begin{derivation}
  1. & $!B \lif (!B \lor !A)$ \\
  2. & $(!B \lif (!B \lor !A)) \lif (!A  \lif (!B \lif (!B \lor !A)))$\\
  3. & $!A  \lif (!B \lif (!B \lor !A))$
\end{derivation}
په ۱مه کرښه !!{sentence} د بديهي اصل $!A \lif (!A \lor !B)$ بڼه
لري، چې د $!A$ او $!B$ ونډې پکې بدلې شوې دي۔ په ۲مه کرښه جمله د
بديهي اصل $!A \lif (!B \lif !A)$ بڼه لري۔ نو دواړه کرښې توجيه شوې
دي۔ ۳مه کرښه د مودوس پوننس له مخې توجيه شوې ده: که په لنډه ئې
$!D$ وليکو، نو ۲مه کرښه د $!C \lif !D$ بڼه لري، چې $!C$ هم
$!B \lif (!B \lor !A)$، يعنې ۱مه کرښه، ده۔

يو سټ $\Gamma$ هله ناسازګار دے چې $\Gamma \Proves \lfalse$۔ د
بديهي اصولو يو بشپړ نظام به دا هم ثابتوي چې د هر~$!A$ دپاره
$\lfalse \lif !A$؛ نو که $\Gamma$ ناسازګار وي، د هر~$!A$ دپاره
$\Gamma \Proves !A$۔

د منطق د بديهي !!{derivation}s نظامونه لومړے ګوټلوب فرېګه په خپل
۱۸۷۹ز اثر \emph{بېګريفشريفت} کښې وړاندې کړل، چې له همدې امله زياتره
د معاصر منطق لومړے اثر ګڼل کېږي۔ بيا الفرېډ نارت وايټ هېډ او برټرېنډ
رسل په \emph{پرېنسيپيا ماتيماتيکا} کښې، او ډيوېډ هيلبرټ او شاګردانو
ئې په ۱۹۲۰مو کلونو کښې بشپړ کړل۔ ځکه ورته زياتره ``د فرېګه نظامونه''
يا ``د هيلبرټ نظامونه'' ويل کېږي۔ دا ډېر انعطاف منونکي دي، ځکه چې د
يو منطق دپاره بديهي نظام زياتره په اسانۍ موندل کېږي۔ دا چې
!!{derivation}s ډېر ساده جوړښت او يوازې يوه يا دوه استنتاجي قاعدې
لري، د هغو \emph{په اړه} د څيزونو ثابتول هم نسبتاً اسانه دي۔ خو په
عمل کښې ئې کارول ډېر ګران دي؛ يعنې ثبوتونه موندل او ليکل ګران دي۔

\end{document}
